


Generated by \texttt{python -m crashguard.rd\_xsim\_doc}. Build record: \texttt{docs/03\_results/XSIM\_BUILD.md}. Scenario code: \texttt{xsim/ns3/xsim\_calib.cc} (ns-3) and \texttt{src/crashguard/rd\_xsim.py} (both sides). Constants: \texttt{docs/03\_results/ASSUMPTIONS.md}.

\textbf{Scenario.} One gateway at the centre of a disc of radius 3,000 m; N vehicles placed uniformly (a new placement and shadowing draw in every run); one spreading factor; 97-byte frames every 30 s on average (periodic with a random phase in ns-3, Poisson arrivals in the custom simulator, a stated difference); one transmission and no retry; unconfirmed or confirmed uplink; no extra random loss. Channel on both sides: path-loss exponent 3 with 25 dB at 1 m, 10 dB vehicle loss, static log-normal shadowing of 6 dB per vehicle, 2 dB fading per packet, 6 dB same-spreading-factor capture, gateway and device sensitivities of the custom simulator (-123, -126, -129 dBm for SF7, SF8, SF9). ns-3 uses its own receiver and interference model, so the comparison tests the two simulators' collision and reception logic, not only the shared constants. N = 1, 10, 50, 100; SF7 and SF9. Samples: 60 hours (custom) and 60 runs (ns-3) per cell.

\textbf{Pre-stated agreement criterion} (\texttt{PREREGISTRATION\_V12.md} section 8): the absolute PDR difference is at most 5 points in every baseline cell.

\subsection*{Results}

{\scriptsize\raggedright
\begin{longtable}{p{0.136\linewidth}p{0.136\linewidth}p{0.136\linewidth}p{0.136\linewidth}p{0.136\linewidth}p{0.136\linewidth}p{0.136\linewidth}}
\toprule
N & SF & Uplink & Custom PDR & ns-3 PDR & Difference (custom minus ns-3) & Within 5 points \\
\midrule
\endhead
1 & 7 & unconfirmed & 0.658 & 0.642 & +0.016 & yes \\
1 & 7 & confirmed & 0.585 & 0.640 & -0.054 & NO \\
1 & 9 & unconfirmed & 0.884 & 0.873 & +0.011 & yes \\
1 & 9 & confirmed & 0.880 & 0.872 & +0.007 & yes \\
10 & 7 & unconfirmed & 0.630 & 0.634 & -0.004 & yes \\
10 & 7 & confirmed & 0.609 & 0.614 & -0.005 & yes \\
10 & 9 & unconfirmed & 0.737 & 0.702 & +0.034 & yes \\
10 & 9 & confirmed & 0.664 & 0.620 & +0.044 & yes \\
50 & 7 & unconfirmed & 0.495 & 0.502 & -0.008 & yes \\
50 & 7 & confirmed & 0.427 & 0.430 & -0.004 & yes \\
50 & 9 & unconfirmed & 0.316 & 0.299 & +0.017 & yes \\
50 & 9 & confirmed & 0.271 & 0.240 & +0.030 & yes \\
100 & 7 & unconfirmed & 0.371 & 0.362 & +0.009 & yes \\
100 & 7 & confirmed & 0.306 & 0.305 & +0.000 & yes \\
100 & 9 & unconfirmed & 0.148 & 0.111 & +0.037 & yes \\
100 & 9 & confirmed & 0.125 & 0.097 & +0.028 & yes \\
\bottomrule
\end{longtable}}

\textbf{Verdict on the pre-stated criterion: NOT MET.} Largest absolute difference 0.054; cells above 5 points: 1 (N=1 SF7 confirmed).

The cell N=1, SF7, confirmed differs by -0.054. With one vehicle the PDR is set by where that vehicle is placed, so 60 placements give a standard error of about 0.090 on the difference (a binomial approximation); the difference is inside that noise. The preregistered criterion is not relaxed on this account.

\subsubsection*{Supplementary: the four single-vehicle cells with 300 placements per side (not part of the criterion; reported next to the main table)}

{\scriptsize\raggedright
\begin{longtable}{p{0.190\linewidth}p{0.190\linewidth}p{0.190\linewidth}p{0.190\linewidth}p{0.190\linewidth}}
\toprule
SF & Uplink & Custom PDR & ns-3 PDR & Difference \\
\midrule
\endhead
7 & unconfirmed & 0.662 & 0.630 & +0.032 \\
7 & confirmed & 0.649 & 0.629 & +0.020 \\
9 & unconfirmed & 0.883 & 0.855 & +0.027 \\
9 & confirmed & 0.907 & 0.856 & +0.050 \\
\bottomrule
\end{longtable}}

\subsection*{Setting (b): both simulators at ns-3's shipped gateway sensitivity (V12 11.3)}

Gateway sensitivity -130, -132.5, -135, -137.5, -140, -142.5 dBm for SF7 to SF12 in both simulators (in ns-3 the end-device table is set to the same values, a stated deviation from the shipped end-device table); same scenarios and sample sizes as setting (a); the local ADR backoff is off. \textbf{Setting (a) above is the setting of every X1.3 verdict.}

{\scriptsize\raggedright
\begin{longtable}{p{0.136\linewidth}p{0.136\linewidth}p{0.136\linewidth}p{0.136\linewidth}p{0.136\linewidth}p{0.136\linewidth}p{0.136\linewidth}}
\toprule
N & SF & Uplink & Custom PDR & ns-3 PDR & Difference (custom minus ns-3) & Within 5 points \\
\midrule
\endhead
1 & 7 & unconfirmed & 0.920 & 0.898 & +0.022 & yes \\
1 & 7 & confirmed & 0.900 & 0.896 & +0.004 & yes \\
1 & 9 & unconfirmed & 0.973 & 0.963 & +0.010 & yes \\
1 & 9 & confirmed & 0.973 & 0.963 & +0.010 & yes \\
10 & 7 & unconfirmed & 0.851 & 0.836 & +0.015 & yes \\
10 & 7 & confirmed & 0.822 & 0.806 & +0.016 & yes \\
10 & 9 & unconfirmed & 0.794 & 0.764 & +0.031 & yes \\
10 & 9 & confirmed & 0.731 & 0.653 & +0.079 & NO \\
50 & 7 & unconfirmed & 0.643 & 0.624 & +0.019 & yes \\
50 & 7 & confirmed & 0.540 & 0.532 & +0.007 & yes \\
50 & 9 & unconfirmed & 0.332 & 0.303 & +0.029 & yes \\
50 & 9 & confirmed & 0.280 & 0.245 & +0.035 & yes \\
100 & 7 & unconfirmed & 0.457 & 0.422 & +0.035 & yes \\
100 & 7 & confirmed & 0.359 & 0.349 & +0.010 & yes \\
100 & 9 & unconfirmed & 0.151 & 0.108 & +0.043 & yes \\
100 & 9 & confirmed & 0.127 & 0.096 & +0.031 & yes \\
\bottomrule
\end{longtable}}

\textbf{Verdict on the pre-stated criterion under setting (b): NOT MET.} Largest absolute difference 0.079; cells above 5 points: 1 (N=10 SF9 confirmed).

\textbf{Both settings (A3).} Setting (a): criterion NOT MET (largest difference 0.054). Setting (b): NOT MET (largest difference 0.079). The original criterion (setting (a)) was missed by one cell at 5.4 points and is reported as a limitation of the cross-simulator comparison whatever setting (b) shows; no constant of the custom simulator was changed to make either setting pass.

\subsection*{What the calibration found (and what it did not)}
\begin{itemize}
\item \textbf{ns-3's shipped gateway sensitivity is 7 dB better than the value assumed in the custom simulator} (SF7: -130 against -123 dBm; SF9: -135 against -129 dBm; the end-device table is 1 dB different). With the shipped values a quick 5-run check of the same cells gave differences of up to 0.22 (single vehicle, SF7); those runs were not kept after the patch, so a rerun with the unpatched build is NOT\_RUN. The comparison above uses a local patch of the two sensitivity tables in the ns-3 source (recorded in \texttt{XSIM\_BUILD.md}). Which sensitivity is right for the real radio is not settled here; the custom simulator's value is an assumption of this repository.
\item ns-3's local adaptive-data-rate backoff changes the data rate after about 96 uplinks without a downlink; it was switched off in the scenario (\texttt{SetUplinkAdrBit(false)}). Before that fix the delivery of a single vehicle fell with the run length (0.92 at 300 s to 0.60 at 3,600 s), which is how it was found.
\item The ns-3 default transmit power (14 dBm) is kept and the 10 dB vehicle loss is applied in the channel instead.
\item Not run: confirmed-uplink ACK success in ns-3 (the tracker counts MAC receptions only); SF11 and SF12 (the ns-3 ALOHA region has no valid data rate for them); more than one gateway.
\end{itemize}

\subsection*{X1.3 and X1.4}
\begin{itemize}
\item **X1.3, D1 as an ns-3 application and the cross-simulator verdict table: run in session 2 for D1 (not D2, already a negative result): \texttt{docs/03\_results/XSIM\_VERDICTS.md}.** The verdicts of H-D1a and H-D1b differ between the simulators; D1 is simulator-dependent. The calibration criterion above was not met, and the verdicts are stated under that condition.
\item \textbf{X1.4, comparison with published numbers (the ns-3 LoRa simulation paper, ICC 2018; the single-gateway model, IEEE TWC 2021; the real-versus-simulated deployment study, IEEE Sensors J. 2022): NOT\_RUN.} None of the three was opened in session 1.
\end{itemize}


